\documentclass[conference]{IEEEtran}
\usepackage{cite}
\usepackage{amsmath,amssymb,amsfonts}
\usepackage{algorithmic}
\usepackage{graphicx}
\usepackage{textcomp}
\usepackage{xcolor}
\usepackage{booktabs}
\usepackage{hyperref}
\usepackage{multirow}

\def\BibTeX{{\rm B\kern-.05em{\sc i\kern-.025em b}\kern-.08em
    T\kern-.1667em\lower.7ex\hbox{E}\kern-.125emX}}

\begin{document}

\title{EventVault-R: An FPGA-Accelerated Autonomous Escrow Architecture for Real-Time Transient Astronomy on Resource-Constrained Platforms}

\author{\IEEEauthorblockN{1\textsuperscript{st} Author Name}
\IEEEauthorblockA{\textit{Department of Space Science} \\
\textit{University/Institute Name}\\
City, Country \\
email@domain.edu}
}

\maketitle

\begin{abstract}
Next-generation time-domain astronomy demands real-time processing of high-cadence focal-plane data under severe power, bandwidth, and computational constraints. Traditional onboard triage systems employ binary ``keep-or-delete'' policies that irreversibly discard observations before their scientific significance can be assessed, resulting in the permanent loss of early-rising transient photometry. We present EventVault-R, a heterogeneous System-on-Chip (SoC) architecture that resolves this problem through three tightly coupled innovations: (1) a deeply pipelined 3-level 2D Biorthogonal 4.4 Discrete Wavelet Transform (DWT) engine implemented in SystemVerilog, (2) a hardware science guardrail that enforces photometric and astrometric fidelity bounds without software intervention, and (3) a cyclic escrow buffer managed by the ARM Hard Processor System (HPS) that defers irreversible data destruction until transient confirmation. Synthesized on an Intel Cyclone V FPGA (5CSEMA5F31C6), the accelerator achieves a worst-case maximum operating frequency ($F_{max}$) of 88.45~MHz, delivering a peak throughput of 88.45~Megapixels per second while consuming only 3.2\% of available Adaptive Logic Modules (ALMs) and 5.7\% of on-chip block memory. Compared to a software-only baseline executing on the same SoC's ARM Cortex-A9 processor, the FPGA fabric delivers a \textbf{327$\times$ improvement in frame processing throughput}. System-level validation confirms 100\% recovery of pre-trigger astronomical observations with 60\% aggregate bandwidth reduction, demonstrating the viability of autonomous, hardware-accelerated science-data triage for CubeSat and edge-astronomy missions.
\end{abstract}

\begin{IEEEkeywords}
FPGA Hardware Acceleration, Discrete Wavelet Transform, Transient Astronomy, System-on-Chip, CubeSat Onboard Processing, Escrow Architecture
\end{IEEEkeywords}

% ===========================================================================
\section{Introduction}
% ===========================================================================

The discovery of astronomical transients---supernovae, gamma-ray burst optical afterglows, tidal disruption events, and fast radio bursts---has transformed our understanding of high-energy astrophysics. However, the physical mechanisms governing the earliest moments of these events remain poorly constrained, precisely because the critical ``early-rising'' photometric data is almost always lost before ground-based confirmation can occur.

Consider a CubeSat constellation monitoring a wide field of sky at 30-second cadence. A Type~Ia supernova begins to brighten at time $T_0$. For the first 10--15 frames, the source's flux increase is buried within the Poisson noise floor of the detector, registering a signal-to-noise ratio (SNR) well below any reasonable detection threshold. By the time the transient crosses the $5\sigma$ detection boundary at time $T_{\text{trig}}$, perhaps 12 minutes later, the onboard triage algorithm has long since purged the pre-trigger observations to free memory for new acquisitions. The most scientifically valuable data---the shock breakout signature, the initial rise curve, the earliest color evolution---is gone forever.

This is not a theoretical concern. Current and planned missions including the Vera C. Rubin Observatory's Legacy Survey of Space and Time (LSST) and ESA's PLATO mission will generate data volumes exceeding 20~TB per night, with CubeSat downlinks often limited to $<$10~kbps. The fundamental tension between \textit{observation rate} and \textit{downlink capacity} forces onboard systems to make irreversible triage decisions under deep uncertainty.

\subsection{The Limitations of Existing Approaches}

Existing onboard data handling architectures fall into two broad categories, both of which are fundamentally inadequate for transient recovery:

\textbf{Store-and-Forward:} The instrument captures all frames and queues them for eventual downlink. This approach preserves all data but is immediately infeasible at high cadences. A $64 \times 64$ pixel, 16-bit focal plane operating at 1~kHz generates 8.19~MB/s of raw data---orders of magnitude beyond any realistic CubeSat downlink budget. The onboard storage fills within minutes.

\textbf{Threshold Triage:} A simple SNR or flux-difference filter is applied to each incoming frame in real time. Frames exceeding the threshold are retained; all others are immediately and irrecoverably deleted. While bandwidth-efficient, this approach is scientifically destructive: by definition, it cannot retain data from events that have not yet crossed the detection threshold.

Neither architecture can simultaneously achieve high bandwidth efficiency \textit{and} retroactive recovery of sub-threshold observations. EventVault-R is designed to break this impasse.

% ===========================================================================
\section{Proposed Solution: The Escrow Architecture}
% ===========================================================================

EventVault-R introduces a fundamentally different paradigm: \textit{deferred data destruction}. Rather than making an immediate binary keep-or-delete decision, the system holds data in a temporary ``escrow'' state, buying time for the transient to reveal itself while bounding the memory and bandwidth cost of this speculative retention.

\subsection{Progressive Wavelet Decomposition}

The key insight enabling this architecture is that raw pixel data is not the optimal representation for speculative retention. Instead, EventVault-R decomposes each incoming frame using a 3-level 2D Discrete Wavelet Transform (DWT) with the Biorthogonal 4.4 (bior4.4) wavelet basis. This produces four distinct resolution layers:

\begin{itemize}
    \item $L_0$ (LL subband): The coarse approximation containing $>$95\% of the image energy. This layer alone is sufficient for source detection and basic photometry.
    \item $H_1$ (Level 1 details): LH, HL, HH subbands capturing fine-scale edges.
    \item $H_2$ (Level 2 details): Medium-scale structural information.
    \item $H_3$ (Level 3 details): Large-scale residuals.
\end{itemize}

The $L_0$ base layer is always retained and transmitted at baseline priority. The detail layers $H_1$--$H_3$ are pushed into a finite circular escrow buffer, where they remain available for a configurable time horizon $\Delta T_{\text{pre}}$. If no trigger occurs within this window, the oldest detail layers are silently purged to make room for new observations. If a trigger \textit{does} occur, the system ``promotes'' all escrowed detail layers within the $[T_{\text{trig}} - \Delta T_{\text{pre}},\ T_{\text{trig}}]$ window to permanent storage for priority downlink.

\subsection{Science Guardrail}

A critical safeguard prevents the system from discarding detail layers that would compromise the scientific validity of the retained base layer. Before any purge decision, the Science Guardrail evaluates two strict conditions against a known source at position $(x_{\text{ref}}, y_{\text{ref}})$:

\begin{equation}
    \frac{|\hat{F} - F_{\text{ref}}|}{F_{\text{ref}}} \leq \epsilon_F \quad (\text{default } 0.5\%)
    \label{eq:flux_guard}
\end{equation}
\begin{equation}
    \sqrt{(\hat{x} - x_{\text{ref}})^2 + (\hat{y} - y_{\text{ref}})^2} \leq \epsilon_x \quad (\text{default } 0.1 \text{ px})
    \label{eq:centroid_guard}
\end{equation}

where $\hat{F}$ is the measured photometric flux and $(\hat{x}, \hat{y})$ is the measured centroid position from the reconstructed image. If either bound is violated, the detail layer is \textit{locked} and cannot be purged, guaranteeing zero uncontrolled science data loss.

\subsection{Why Hardware Acceleration Is Necessary}

The mathematical operations involved in a single frame's 3-level 2D bior4.4 DWT decomposition are substantial. For a $64 \times 64$ input, the bior4.4 wavelet requires 9-tap (decomposition low-pass) and 7-tap (decomposition high-pass) filters applied across rows and columns at each of three decomposition levels. On a general-purpose ARM Cortex-A9 processor (the HPS embedded in the Cyclone V SoC), this computation is executed sequentially: each multiply-accumulate (MAC) operation occupies the single ALU pipeline for multiple clock cycles, and the entire $64 \times 64$ frame requires approximately $15$~ms to process.

At this rate, the processor is limited to $\approx$66 frames per second (FPS)---adequate for low-cadence survey modes but wholly insufficient for high-speed event characterization, where frame rates of 1--10~kHz are routine. Furthermore, the processor's power consumption during sustained DWT computation far exceeds the $<$2~W power budgets typical of CubeSat payloads.

An FPGA implementation, by contrast, exploits \textit{spatial parallelism}: the filter taps are unrolled into dedicated hardware multipliers (DSP blocks), and the row and column transforms are orchestrated by a deterministic finite state machine operating at wire speed. The result is a fully pipelined architecture that produces one output coefficient per clock cycle, achieving throughput that scales linearly with clock frequency rather than being bottlenecked by instruction fetch and decode overhead.

% ===========================================================================
\section{Methodology: Hardware Implementation}
% ===========================================================================

The EventVault-R hardware accelerator was implemented in SystemVerilog and targets the Intel Cyclone V SoC (5CSEMA5F31C6) on the Terasic DE1-SoC development board. The design follows a three-phase development methodology: (A) Python reference model, (B) embedded C++ port with Q16.16 fixed-point arithmetic, and (C) RTL hardware acceleration with formal verification against the software golden vectors.

\subsection{DWT Engine Architecture}

The core DWT engine consists of three hierarchical modules:

\textbf{dwt\_1d\_bior.sv} implements a single 1D bior4.4 forward wavelet transform. Input samples are written into a synchronous buffer annotated with \texttt{(* ramstyle = "M10K" *)} to force inference into dedicated block RAM. A finite state machine applies the 9/7 filter bank using symmetric boundary extension and Q16.16 fixed-point multiply-accumulate operations. The design was carefully pipelined to account for the 1-cycle read latency inherent in M10K synchronous reads: the address is presented on cycle $N$ and the data is consumed on cycle $N+1$.

\textbf{dwt\_2d\_level.sv} orchestrates a complete single-level 2D DWT by applying the 1D engine first across all rows, storing the intermediate L and H results in M10K block RAMs, and then applying the 1D engine across columns to produce the four subbands (LL, LH, HL, HH). Six large intermediate arrays (\texttt{L\_rows}, \texttt{H\_rows}, \texttt{res\_LL}, \texttt{res\_LH}, \texttt{res\_HL}, \texttt{res\_HH}) are all explicitly mapped to M10K blocks using linear address counters to avoid multiplier-based indexing.

\textbf{dwt\_2d\_top.sv} chains three instances of the 2D level module to perform the full 3-level decomposition. The LL output of each level is buffered in an M10K-backed array and sequentially fed into the next level. The final outputs are the $L_0$ (Level 3 LL), $H_1$ (Level 1 details), $H_2$ (Level 2 details), and $H_3$ (Level 3 details) subbands.

\subsection{Science Guardrail Hardware}

The hardware guardrail (\texttt{science\_guardrail.sv}) operates on the reconstructed pixel stream using a 5$\times$5 sliding window backed by a 5-line circular buffer. It computes the photometric flux (sum of pixel intensities) and weighted centroid displacements ($\Delta x$, $\Delta y$) in a 4-stage pipeline:

\begin{enumerate}
    \item \textbf{Stage 1:} Per-row flux accumulation and $dx$-weighted sums.
    \item \textbf{Stage 2:} Cross-row summation for total flux, $\Sigma dx$, and $\Sigma dy$.
    \item \textbf{Stage 3:} Reference multiplication to compute bounds ($\epsilon_F \cdot F_{\text{ref}}$ and $\epsilon_x \cdot F$) in Q15 fixed-point.
    \item \textbf{Stage 4:} Comparison and \texttt{is\_safe} assertion.
\end{enumerate}

This division-free formulation avoids expensive hardware dividers entirely by reformulating the centroid error as $|\Sigma dx| \leq \epsilon_x \cdot F$ (comparing against the product rather than dividing by $F$).

\subsection{Memory Architecture Optimization}

An initial synthesis attempt mapped all intermediate DWT arrays to logic registers, consuming 6,422 LABs---exceeding the Cyclone V's 3,207 LAB capacity by nearly 2$\times$. The root cause was identified as combinational (asynchronous) reads from array elements, which prevents Quartus from inferring M10K block RAMs.

The solution required three coordinated changes across all RTL modules:
\begin{enumerate}
    \item Annotating all large arrays with \texttt{(* ramstyle = "M10K" *)}.
    \item Replacing all combinational reads with synchronous \texttt{always\_ff} read pipelines.
    \item Replacing multiplier-based array indexing (e.g., \texttt{row * MAX\_COLS + col}) with linear address counters to satisfy the single-address-port constraint of M10K blocks.
\end{enumerate}

This optimization reduced LAB consumption by \textbf{84\%} (from 6,422 to 1,032) while shifting the storage burden to the abundant M10K block RAM resources.

\subsection{Fixed-Point Arithmetic: Q16.16}

All arithmetic within the FPGA fabric uses a 32-bit Q16.16 fixed-point representation, providing 16 bits of integer range and 16 bits of fractional precision. The bior4.4 filter coefficients (which have magnitudes $<$2.0) are pre-scaled to Q16.16 constants at synthesis time. Intermediate MAC accumulations use 64-bit products, which are right-shifted by 16 bits before truncation to maintain precision across the filter bank.

A dedicated C++ reference model (\texttt{dwt\_fixed.cpp}) mirrors the exact Q16.16 arithmetic of the hardware, enabling bit-accurate verification of the RTL against the software golden vectors.

% ===========================================================================
\section{Results and Analysis}
% ===========================================================================

\subsection{FPGA Resource Utilization}

Table~\ref{tab:utilization} summarizes the post-place-and-route resource consumption reported by the Quartus Prime 25.1 Fitter for the complete EventVault-R accelerator (DWT engine + science guardrail + AXI wrapper).

\begin{table}[ht]
\centering
\caption{Cyclone V FPGA Resource Utilization (Post-Fit)}
\label{tab:utilization}
\begin{tabular}{@{}lccc@{}}
\toprule
\textbf{Resource} & \textbf{Used} & \textbf{Available} & \textbf{Utilization} \\ \midrule
Adaptive Logic Modules (ALMs) & 1,032 & 32,070 & 3.2\% \\
Dedicated Logic Registers & 1,076 & 64,140 & 1.7\% \\
Block Memory Bits (M10K) & 234,656 & 4,065,280 & 5.8\% \\
M10K RAM Blocks & 44 & 397 & 11.1\% \\
DSP Blocks (18$\times$18 Multipliers) & 15 & 87 & 17.2\% \\
PLLs & 0 & 6 & 0\% \\ \bottomrule
\end{tabular}
\end{table}

The accelerator occupies a remarkably small fraction of the target device. The 3.2\% ALM utilization leaves $>$96\% of the logic fabric available for additional mission payloads (e.g., communication controllers, attitude determination, or additional sensor pipelines). The 11.1\% M10K utilization confirms that the memory optimization strategy successfully shifted the storage burden from scarce logic registers to abundant block RAM.

\subsection{Timing Performance}

Static Timing Analysis (STA) was performed under both the Slow 1100mV 85$^\circ$C and Slow 1100mV 0$^\circ$C corner models. The results are summarized in Table~\ref{tab:timing}.

\begin{table}[ht]
\centering
\caption{Static Timing Analysis Results}
\label{tab:timing}
\begin{tabular}{@{}lcc@{}}
\toprule
\textbf{Corner Model} & \textbf{$F_{max}$} & \textbf{Margin over 50~MHz Target} \\ \midrule
Slow 1100mV, 85$^\circ$C & 90.30 MHz & +80.6\% \\
Slow 1100mV, 0$^\circ$C & 88.45 MHz & +76.9\% \\ \bottomrule
\end{tabular}
\end{table}

Even under the most pessimistic timing corner (88.45~MHz), the design exceeds the target 50~MHz board clock by 76.9\%, providing substantial timing margin for reliability. At the achievable $F_{max}$, the accelerator delivers a peak throughput of \textbf{88.45 Megapixels per second}.

\subsection{Hardware vs. Software Performance Comparison}

To quantify the acceleration benefit, the identical 3-level 2D bior4.4 DWT algorithm was executed on the DE1-SoC's integrated ARM Cortex-A9 dual-core processor (800~MHz) using optimized C++17 with \texttt{-O2} compiler flags. Table~\ref{tab:comparison} presents the head-to-head comparison for a $64 \times 64$ pixel region of interest.

\begin{table}[ht]
\centering
\caption{Software (ARM Cortex-A9) vs. Hardware (FPGA) Performance}
\label{tab:comparison}
\begin{tabular}{@{}p{3.2cm}cc@{}}
\toprule
\textbf{Metric} & \textbf{ARM CPU} & \textbf{FPGA Accelerator} \\ \midrule
Clock Frequency & 800 MHz & 88.45 MHz \\
Processing Latency & $\approx$15.0 ms & $\approx$0.046 ms \\
Max Frame Rate & 66 FPS & $>$16,000 FPS \\
Peak Throughput & 0.27 MP/s & \textbf{88.45 MP/s} \\
Speedup Factor & $1\times$ (baseline) & \textbf{327$\times$} \\
Processing Model & Sequential & Pipelined Parallel \\
Power Profile & Active CPU core & Dedicated fabric \\ \bottomrule
\end{tabular}
\end{table}

The FPGA accelerator achieves a \textbf{327$\times$ throughput improvement} over the embedded processor despite operating at a clock frequency \textit{9$\times$ lower} than the ARM core. This counter-intuitive result underscores the fundamental advantage of spatial parallelism: while the ARM processor executes one multiply-accumulate per cycle through its general-purpose ALU, the FPGA fabric instantiates 15 dedicated DSP multipliers operating simultaneously, with results flowing through a deterministic pipeline without instruction-fetch overhead, branch prediction penalties, or cache miss stalls.

For a high-speed astronomical camera operating at 1,000 FPS with $64 \times 64$ pixel regions, the ARM processor would require 15~seconds of computation per second of observation---an obvious impossibility. The FPGA accelerator, by contrast, completes the same workload in 46~milliseconds, leaving 954~ms of margin per second for housekeeping, communication, and power management.

\subsection{Mathematical Fidelity of Fixed-Point Arithmetic}

The transition from double-precision floating-point (64-bit IEEE 754) to Q16.16 fixed-point (32-bit) arithmetic introduces quantization error. Table~\ref{tab:fidelity} summarizes the distortion measured by comparing the C++ fixed-point pipeline against the Python double-precision reference model across all test vectors.

\begin{table}[ht]
\centering
\caption{Fixed-Point (Q16.16) vs. Floating-Point Fidelity}
\label{tab:fidelity}
\begin{tabular}{@{}lcc@{}}
\toprule
\textbf{Metric} & \textbf{Float $\rightarrow$ Fixed Error} & \textbf{Tolerance} \\ \midrule
RMSE (pixel intensity) & $<$0.1 units & --- \\
Max Absolute Error & $<$0.5 units & --- \\
Photometric Flux Error ($\Delta F$) & $<$0.5\% & $\leq$0.5\% \\
Astrometric Centroid Error ($\Delta x$) & $<$0.1 px & $\leq$0.1 px \\
Science Guardrail Verdict & \textbf{PASS} & --- \\ \bottomrule
\end{tabular}
\end{table}

The Q16.16 representation provides 16 bits of fractional precision ($\approx$1.5$\times 10^{-5}$ resolution), which is more than sufficient to preserve the bior4.4 filter coefficients and the resulting subband energies without exceeding the science guardrail bounds. This result validates the architectural decision to avoid floating-point hardware entirely, eliminating the need for power-hungry FPU IP cores and enabling the compact 3.2\% ALM footprint.

\subsection{System-Level Efficacy: Bandwidth Reduction}

The overarching purpose of EventVault-R is not merely fast DWT computation, but autonomous, lossless-where-it-matters bandwidth reduction. To validate this system-level behavior, a 50-frame synthetic observation sequence was constructed with the following properties:

\begin{itemize}
    \item Frames $t=0$ to $t=49$ contain a static star field with Poisson noise.
    \item A transient source begins brightening at $t=10$ but remains below the $5\sigma$ detection threshold until $t=25$.
    \item The detection trigger fires at $t=25$.
\end{itemize}

Table~\ref{tab:escrow} summarizes the system's response.

\begin{table}[ht]
\centering
\caption{Escrow Promotion and Bandwidth Reduction}
\label{tab:escrow}
\begin{tabular}{@{}lcc@{}}
\toprule
\textbf{Metric} & \textbf{Naive Triage} & \textbf{EventVault-R} \\ \midrule
Pre-trigger frames recovered & 0 / 15 & \textbf{15 / 15} \\
Pre-trigger recovery rate & 0\% & \textbf{100\%} \\
Frames transmitted (of 50) & 25 & 20 \\
Bandwidth reduction & 50\% & \textbf{60\%} \\
Uncontrolled data loss events & 15 & \textbf{0} \\
Buffer overflows & 0 & \textbf{0} \\ \bottomrule
\end{tabular}
\end{table}

The naive threshold triage baseline successfully detected the transient at $t=25$ and retained frames $t=25$ through $t=49$, achieving 50\% bandwidth reduction. However, it \textit{permanently lost} all 15 pre-trigger frames ($t=10$ through $t=24$), destroying the early-rising photometric data that is most scientifically valuable.

EventVault-R, by contrast, retained the detail layers of all 50 frames in its circular escrow buffer. Upon trigger at $t=25$, the promotion controller swept the buffer and recovered \textbf{all 15 pre-trigger frames} with full multi-resolution fidelity. The system then safely purged the 30 non-anomalous baseline frames ($t=0$ through $t=9$ and $t=36$ through $t=49$), achieving a superior \textbf{60\% bandwidth reduction} while preserving 100\% of the scientifically critical data.

% ===========================================================================
\section{Conclusion}
% ===========================================================================

EventVault-R demonstrates that the tension between real-time data rates and constrained downlink budgets can be resolved through a heterogeneous SoC architecture that combines dedicated FPGA acceleration with intelligent software-managed escrow buffering. The key contributions of this work are:

\begin{enumerate}
    \item A deeply pipelined 3-level 2D bior4.4 DWT engine that achieves 88.45~MP/s throughput using only 3.2\% of a Cyclone V FPGA's logic resources---a \textbf{327$\times$ speedup} over the embedded ARM processor.
    \item A division-free hardware science guardrail that enforces photometric ($\Delta F \leq 0.5\%$) and astrometric ($\Delta x \leq 0.1$~px) fidelity bounds at wire speed.
    \item A cyclic escrow architecture that achieves \textbf{100\% recovery} of pre-trigger astronomical observations while reducing aggregate bandwidth by 60\%.
    \item A complete verification methodology spanning Python golden vectors, C++ Q16.16 bit-accurate models, Icarus Verilog RTL simulation, and Quartus FPGA synthesis.
\end{enumerate}

By deferring irreversible data destruction until transient confirmation, EventVault-R ensures that the most scientifically valuable observations---the ones that exist \textit{before} anyone knows they matter---are never lost.

\begin{thebibliography}{1}
\bibitem{mallat1989} S.~Mallat, ``A theory for multiresolution signal decomposition: the wavelet representation,'' \textit{IEEE Trans. Pattern Anal. Mach. Intell.}, vol.~11, no.~7, pp.~674--693, Jul. 1989.
\bibitem{cohen1992} A.~Cohen, I.~Daubechies, and J.-C.~Feauveau, ``Biorthogonal bases of compactly supported wavelets,'' \textit{Comm. Pure Appl. Math.}, vol.~45, no.~5, pp.~485--560, 1992.
\bibitem{altera_m10k} Intel Corporation, ``Cyclone V Device Handbook,'' vol.~1, 2018. [Online]. Available: https://www.intel.com/content/www/us/en/docs/programmable/683694/
\end{thebibliography}

\end{document}
